\chapter{Introducción}

La representación de fenómenos atmosféricos en tiempo real constituye uno de los problemas clásicos dentro del área de gráficos por computadora. Efectos como la lluvia, la nieve, la niebla o el humo pueden involucrar miles de elementos visibles de manera simultánea, lo que representa un desafío importante para el rendimiento de una aplicación gráfica. Una implementación basada en geometría tridimensional tradicional incrementaría considerablemente el número de polígonos procesados por la GPU, reduciendo la velocidad de renderizado y limitando la capacidad de mantener una tasa de cuadros estable.

Una de las técnicas más utilizadas para resolver este problema consiste en el empleo de \textit{billboards}. Un billboard es un plano bidimensional texturizado que modifica automáticamente su orientación para permanecer siempre visible desde la posición de la cámara. Esta técnica permite representar objetos visualmente complejos utilizando únicamente un cuadrilátero, disminuyendo significativamente la carga de procesamiento sin afectar de manera apreciable la percepción visual del usuario.

Además del uso de billboards, los motores gráficos modernos emplean sistemas de partículas para simular fenómenos naturales. Un sistema de partículas administra la creación, actualización y reutilización de una gran cantidad de elementos independientes, permitiendo representar efectos dinámicos como lluvia, fuego, humo, polvo o explosiones. Debido a que cada partícula posee un ciclo de vida propio, el sistema puede reciclar continuamente los recursos gráficos disponibles, logrando un alto rendimiento incluso cuando se representan miles de partículas de manera simultánea.

En este proyecto se desarrolló una aplicación interactiva utilizando OpenGL 3.3 y el lenguaje C++, cuyo objetivo principal consiste en simular una tormenta de lluvia en tiempo real mediante la utilización de billboards y un sistema de partículas. La escena incorpora además un terreno texturizado, árboles implementados mediante billboards, un sistema básico de clima con relámpagos, sonido ambiental de lluvia y una interfaz gráfica desarrollada con Dear ImGui que permite modificar diversos parámetros de ejecución en tiempo real.

Durante el desarrollo se adoptó una arquitectura modular basada en componentes independientes, separando claramente las responsabilidades de renderizado, administración de escenas, recursos gráficos, partículas, cámara, interfaz gráfica, sistema de audio y efectos climáticos. Esta organización facilita la reutilización del código, mejora la mantenibilidad del proyecto y permite incorporar nuevas funcionalidades sin modificar significativamente los módulos existentes.

Finalmente, el proyecto demuestra que la combinación de billboards, sistemas de partículas y una adecuada administración de recursos constituye una solución eficiente para la representación de fenómenos atmosféricos en aplicaciones gráficas en tiempo real, manteniendo un equilibrio entre calidad visual y rendimiento computacional.